Entre los 91 pares de perfiles, el número de dominios con un estado interpretable y consistente en ambos oscila entre 45 y 50 en registros A, pero se reduce a un rango de 10 a 14 en registros AAAA debido a que solo 11 de los 50 dominios publicaban registros IPv6 en la línea base sin filtro. Por ejemplo, Cloudflare Estándar y Familiar coinciden en 46 de 50 dominios en A (acuerdo 0,92) y en 10 de 11 en AAAA (acuerdo 0,91), excluyendo 39 dominios en este último caso; ambos acuerdos responden así a coberturas muy distintas.

En la categoría documental de contenido adulto, la regla automática común registró 12 señales en A y 12 en AAAA para Cloudflare Familiar, y 12 señales en AAAA para AdGuard Familiar. La inspección de los mensajes completos por proveedor explica por qué los recuentos uniformes por sí solos no permiten ordenar a los servicios:

\begin{itemize}
\item En A, AdGuard Familiar devolvió \texttt{94.140.14.35} en los cuatro dominios y las tres rondas, dirección identificada por el equipo de AdGuard como su página de bloqueo \cite{adguardblock}.
\item CleanBrowsing Adulto y Familiar devolvieron \texttt{NXDOMAIN} en los cuatro dominios, tanto en A como en AAAA, en concordancia con su guía oficial de verificación \cite{cleanbrowsingverify}. Los controles positivos en A distinguen estos casos de una caída general; en AAAA, la ausencia de direcciones en la referencia limita el poder discriminante de la comparación.
\item Control D Familiar devolvió \texttt{0.0.0.0} y \texttt{::} sin código EDE en los cuatro dominios, formato consistente con su documentación de bloqueo \cite{controldblock}. Bajo la regla común sin huella específica para ese operador, estas respuestas quedaron conservadas como anomalías.
\item OpenDNS FamilyShield devolvió \texttt{146.112.61.106} en A y su dirección IPv4 mapeada en IPv6 (\texttt{::ffff:146.112.61.106}) en AAAA, huellas documentadas por Cisco como página de bloqueo por categoría y confirmadas por separado frente a su dominio de prueba y al perfil estándar \cite{ciscoblock,ciscotest}.
\end{itemize}

Esta inspección demuestra que una dirección IPv4 o IPv6 válida devuelta por DNS puede corresponder a una página de bloqueo y que la comparación rigurosa entre servicios PDNS exige combinar reglas comunes con huellas verificadas por operador.
